% docspec-cas frontmatter（無作者：只給標題與副標）
\let\WriteBookmarks\relax
\def\floatpagefraction{.9}

% ---------- #5 移除空白 ABSTRACT 區 ----------
% upstream cas-sc 的 \maketitle→\MaketitleBox 無條件跑 \begin{Abstract}...\end{Abstract}，
% 即使沒有摘要也會印出「ABSTRACT」標題＋分隔線。本文件無摘要、首段即正文導言，
% 故把 Abstract 環境改成「吞掉內容、不輸出任何東西」，並把它前後的 \dashrule
% （upstream cas-common 中僅用於摘要區）改成 no-op。標題之後直接進正文流。
\RenewDocumentEnvironment{Abstract}{o}{%
  \setbox0=\vbox\bgroup%
}{%
  \egroup%
}
\RenewDocumentCommand{\dashrule}{O{.4pt} m m}{}

% ---------- #4 移除「Preprint submitted to Elsevier」頁尾 ----------
% upstream cas-sc 兩個頁面樣式（ps@first 第一頁、ps@cas 其餘）的頁尾都呼叫
% \__first_footerline:，其中含 short authors + 斜體「Preprint submitted to Elsevier」。
% 把它清空即可拿掉該行；右側「Page X of Y」由 \__first_foot:/\__cas_foot: 另行輸出，
% 不受影響、仍保留。
\ExplSyntaxOn
\cs_set:Npn \__first_footerline: { }
\ExplSyntaxOff

% 無作者→無 ORCID，但 upstream cas-sc 的 \printFirstPageNotes 仍無條件呼叫 \printorcid，
% 在標題註腳區留下空的「ORCID(s):」標籤（與空 ABSTRACT 同類殘留）。清成 no-op。
\RenewDocumentCommand{\printorcid}{}{}

% ---------- #5 移除標題下方 affiliation / nonumnote 殘留區 ----------
% 標題下方「開發者實作指引（依台中港車隊 ICD v1.1.0）」原由 \nonumnote 注入，
% 經 \printnonumnotes 在標題註腳區印出，且註腳區頂帶一條 \rule 分隔線（upstream cas-common
% 的 \printFirstPageNotes footnote 文字框）。本文件不需此 affiliation 區：
%   1) 不再呼叫 \nonumnote（見下方標題段，已移除該行）；
%   2) \printnonumnotes 清成 no-op（雙保險：即使殘留 note 也不輸出）。
% 標題之後直接進正文，無多餘橫線/空塊。
\renewcommand{\printnonumnotes}{}

% 標題由 docspec export 在 build 時注入：__DOCSPEC_TITLE__ 取自快照首個 H1
% （無 H1 → 退回文章名）。shorttitle 用同一字串。
\shorttitle{__DOCSPEC_TITLE__}
\shortauthors{docspec}

% 標題置中：upstream cas-sc 的 stm/title「mode=title」meta 預設 align=\raggedright（靠左）。
% \title 的 O{} 選項會在套用 mode meta 之後再跑一次 \keys_set，故在此追加
% align=\centering 即覆蓋對齊，讓主標題水平置中（其餘 size/shape/weight 不動）。
\title[mode=title,align=\centering]{__DOCSPEC_TITLE__}
% #5 原 \nonumnote{開發者實作指引（依台中港車隊 ICD v1.1.0）} 已移除：
% 該行會在標題下方印出帶橫線的 affiliation/note 殘留區。標題後直接進正文。

\maketitle

% ★字級根本修正（C4）：upstream cas-sc 在 \maketitle 期間把 NFSS 字級巨集（\normalsize/\small/...）
% 重設回 class 預設（≈12pt），故 preamble 裡的 \renewcommand 全失效（內文/表格/行內 code
% 都吃舊小字）。解法＝在 \maketitle **之後**（cas 重設完）才重定義字級巨集，於是內文、表格
% (\fontsize 或 \normalsize)、行內 code(\small) 全部一次吃到目標字級。baseline 為自然值，行距
% 由 preamble 的 \linespread 放大（勿預乘）。★此字級階梯由 docspec 從 base_size 旋鈕確定性
% 編出、在 build 時注入（compile_postmaketitle_fonts）；預設 base_size＝14.5 與舊硬寫值
% byte-identical。base_size 旋鈕因此真的控制內文字級。
__DOCSPEC_FONTSIZES__
